% =====================================================================
% Skeleton E — 5 Stage Horizontal + Multi-Parallel Input + Multi Viz
% =====================================================================
% USAGE: Adapted from a GOLD STANDARD figure (Multi-modal Medical Fusion).
% Sub-agent: copy ENTIRE file as figure.tex, change ONLY content.
%
% LAYOUT MODEL (extension of skeleton B for fusion-style architectures):
%   ┌─────┐    ┌─────┐    ┌─────────┐    ┌─────┐    ┌──────────┐
%   │Mod 1│───▶│Proj │───▶│         │───▶│Head │───▶│ Eval/Vis │
%   │Mod 2│───▶│Proj │───▶│  Hero   │───▶│Head │───▶│   +      │
%   │Mod 3│───▶│Proj │───▶│ Cross-  │───▶│Head │───▶│  Explain │
%   └─────┘    └─────┘    │   Attn  │    └─────┘    │ (SHAP/   │
%                          │  + Viz  │                │  GradCAM)│
%                          └─────────┘                └──────────┘
%   Bottom: multi-method bar chart (compares fusion vs single modality)
%
% WHEN TO USE: Multi-modal architectures (vision+text+audio), fusion
% models, explainable AI with multiple visualization channels.
%
% CONTAINS: 3 parallel modality inputs + projection + cross-attn hero
% with embedded heatmap + 2 task heads + SHAP/GPT panels + bottom bar.
%
% CONSTRAINTS — DO NOT VIOLATE:
%
%   [RENAMING]
%   - 3 modality input names (e.g., CT/MRI/PET) — RENAME to topic
%     modalities (vision+text+audio, RGB+depth+IR, etc.).
%   - Hero "Cross-Attention" — RENAME for non-attention fusion.
%   - 2 task heads — RENAME to topic outputs.
%   - Data annotation box content (dataset name/size) — RENAME.
%
%   [SUB-PANEL CENTERING — these shifts were tuned by hand]
%   - 细分类头 bar chart x-shift = 15.45 (centered at zone x=16.5).
%   - Mini Grad-CAM heatmap x-shift = 20.7 (centered at zone x=21.5).
%     If you move the parent zones, recompute shifts to recenter.
%   - Data annotation box uses anchor=north at (25.5, 15.3).
%     Do NOT change to anchor=north east at (27.0, 15.3) — visually
%     off-center against the zone.
%
%   [MDT FEEDBACK DASHED LINE — critical]
%   - Path MUST be orthogonal U-shape: from MDT down to y=-1.0, then
%     LEFT along y=-1.0 across the figure, then UP to input height,
%     then RIGHT into path input. NEVER use a diagonal — a diagonal
%     line crosses the entire figure body and ruins the layout.
% =====================================================================

\documentclass[tikz,border=30pt]{standalone}
\usepackage{tikz}
\usepackage[fontset=none]{ctex}
\setCJKmainfont{PingFang SC}
\setCJKsansfont{PingFang SC}
\usepackage{amsmath,amssymb}
\usetikzlibrary{shapes, arrows.meta, positioning, fit, backgrounds, calc, shadows, shapes.geometric}

% =====================================================================
% STYLE PRESET: Academic Professional (default)
% =====================================================================
% E's color identity:
%   primary=blue (modality 1), success=green (modality 2),
%   accent=orange (modality 3), hero=purple (cross-attention fusion),
%   warning=red, neutral=grey.
%   4 bar colors for grouped comparison; 4-stop heatmap gradient + bg.
%
% Swap style by replacing PRESET DEFINITION block only.
% See references/style-presets/README.md for token semantics.
% =====================================================================

% ----- BEGIN PRESET DEFINITION -----
\definecolor{c_primary_line}{HTML}{6C8EBF}       \definecolor{c_primary_fill}{HTML}{DAE8FC}
\definecolor{c_hero_line}{HTML}{9673A6}          \definecolor{c_hero_fill}{HTML}{E1D5E7}
\definecolor{c_accent_line}{HTML}{D79B00}        \definecolor{c_accent_fill}{HTML}{FFE6CC}
\definecolor{c_success_line}{HTML}{82B366}       \definecolor{c_success_fill}{HTML}{D5E8D4}
\definecolor{c_warning_line}{HTML}{B85450}       \definecolor{c_warning_fill}{HTML}{F8CECC}
\definecolor{c_neutral_line}{HTML}{666666}       \definecolor{c_neutral_fill}{HTML}{F5F5F5}
\definecolor{c_bar_1}{HTML}{3B82F6}              \definecolor{c_bar_2}{HTML}{22C55E}
\definecolor{c_bar_3}{HTML}{F59E0B}              \definecolor{c_bar_4}{HTML}{EF4444}
\definecolor{c_heat_5}{HTML}{6D28D9}             \definecolor{c_heat_4}{HTML}{3B82F6}
\definecolor{c_heat_3}{HTML}{93C5FD}             \definecolor{c_heat_2}{HTML}{DBEAFE}
\definecolor{c_heat_bg}{HTML}{F0F4FF}
% ----- END PRESET DEFINITION -----

% ----- BEGIN BACKWARDS-COMPAT ALIASES (DO NOT EDIT) -----
\colorlet{drawBlueLine}{c_primary_line}          \colorlet{drawBlueFill}{c_primary_fill}
\colorlet{drawGreenLine}{c_success_line}         \colorlet{drawGreenFill}{c_success_fill}
\colorlet{drawOrangeLine}{c_accent_line}         \colorlet{drawOrangeFill}{c_accent_fill}
\colorlet{drawPurpleLine}{c_hero_line}           \colorlet{drawPurpleFill}{c_hero_fill}
\colorlet{drawRedLine}{c_warning_line}           \colorlet{drawRedFill}{c_warning_fill}
\colorlet{drawGreyLine}{c_neutral_line}          \colorlet{drawGreyFill}{c_neutral_fill}
\colorlet{barBlue}{c_bar_1}                      \colorlet{barGreen}{c_bar_2}
\colorlet{barOrange}{c_bar_3}                    \colorlet{barRed}{c_bar_4}
\colorlet{heatDeep}{c_heat_5}                    \colorlet{heatDark}{c_heat_4}
\colorlet{heatMid}{c_heat_3}                     \colorlet{heatLight}{c_heat_2}
\colorlet{heatBg}{c_heat_bg}
% ----- END BACKWARDS-COMPAT ALIASES -----

\pgfdeclarelayer{background}
\pgfsetlayers{background,main}

\begin{document}
\begin{tikzpicture}[
    every node/.style={font=\small},
    % === Box styles ===
    base_box/.style={rectangle, rounded corners=3pt, align=center,
        minimum height=0.9cm, minimum width=2.6cm,
        inner sep=8pt, drop shadow={opacity=0.12}, thick},
    blue_node/.style={base_box, fill=drawBlueFill, draw=drawBlueLine},
    green_node/.style={base_box, fill=drawGreenFill, draw=drawGreenLine},
    orange_node/.style={base_box, fill=drawOrangeFill, draw=drawOrangeLine},
    purple_node/.style={base_box, fill=drawPurpleFill, draw=drawPurpleLine},
    red_node/.style={base_box, fill=drawRedFill, draw=drawRedLine, font=\small\bfseries},
    grey_node/.style={base_box, fill=drawGreyFill, draw=drawGreyLine},
    small_box/.style={rectangle, rounded corners=2pt, align=center,
        minimum height=0.55cm, minimum width=1.4cm,
        inner sep=4pt, thick, font=\footnotesize},
    proj_box/.style={small_box, fill=drawGreyFill!80, draw=drawGreyLine,
        minimum width=1.0cm, minimum height=1.8cm},
    % === Arrow styles ===
    arrow/.style={-{Stealth[scale=0.9]}, thick, color=black!70,
        shorten >=2pt, shorten <=1pt},
    data_arrow/.style={-{Stealth[scale=1.1]}, very thick, color=drawOrangeLine,
        shorten >=2pt, shorten <=1pt},
    feed_arrow/.style={-{Stealth[scale=0.9]}, thick, dashed, color=drawRedLine,
        shorten >=2pt, shorten <=1pt},
    thin_arrow/.style={-{Stealth[scale=0.7]}, thick, color=drawGreyLine},
    tag/.style={font=\footnotesize, inner sep=2pt, rounded corners=1pt},
]

% ================================================================
%  STAGE 1: MULTIMODAL FEATURE EXTRACTION (x=0~5)
% ================================================================
% --- CT pathway (top) ---
\node[blue_node, minimum width=3.0cm, minimum height=2.0cm] (ct) at (1.5, 14)
    {\textbf{CT 影像}\\[1pt]{\footnotesize 3D ResNet-50}\\[-1pt]{\scriptsize 2048维}};
% 3D pseudo-stereo effect on CT box
\fill[drawBlueFill!40, draw=drawBlueLine!80, thick]
    ([xshift=4pt,yshift=4pt]ct.north east) -- (ct.north east)
    -- (ct.south east) -- ([xshift=4pt,yshift=4pt]ct.south east) -- cycle;
\fill[drawBlueFill!30, draw=drawBlueLine!80, thick]
    ([xshift=4pt,yshift=4pt]ct.north west) -- ([xshift=4pt,yshift=4pt]ct.north east)
    -- (ct.north east) -- (ct.north west) -- cycle;

% --- Pathology pathway (middle) ---
\node[green_node, minimum width=3.0cm, minimum height=2.0cm] (path) at (1.5, 10)
    {\textbf{病理切片}\\[1pt]{\footnotesize ViT-B/16}\\[-1pt]{\scriptsize 768维}};

% --- EHR pathway (bottom) ---
\node[orange_node, minimum width=3.0cm, minimum height=2.0cm] (ehr) at (1.5, 6)
    {\textbf{电子病历}\\[1pt]{\footnotesize 临床 BERT}\\[-1pt]{\scriptsize 768维}};

% --- Projection layers ---
\node[proj_box] (projCT) at (5, 14) {\scriptsize 投影\\[-1pt]\scriptsize 512维};
\node[proj_box] (projPath) at (5, 10) {\scriptsize 投影\\[-1pt]\scriptsize 512维};
\node[proj_box] (projEHR) at (5, 6) {\scriptsize 投影\\[-1pt]\scriptsize 512维};

% ================================================================
%  STAGE 2: CROSS-MODAL TRANSFORMER FUSION (x=7~14)
% ================================================================
% CMT large box (core module, 2x size)
\node[red_node, minimum width=6.5cm, minimum height=9.5cm,
      line width=1.5pt] (cmt) at (10.5, 10) {};
\node[font=\small\bfseries, text=drawRedLine!80!black] at (10.5, 14.0)
    {跨模态 Transformer (CMT)};
\node[font=\footnotesize, text=drawRedLine!60!black] at (10.5, 13.4)
    {6层交叉注意力};

% Inside CMT: Cross-attention diagram (simplified)
\node[small_box, fill=drawBlueFill!50, draw=drawBlueLine!70,
      minimum width=1.8cm] (qry) at (8.8, 12.0) {\scriptsize Query (CT)};
\node[small_box, fill=drawGreenFill!50, draw=drawGreenLine!70,
      minimum width=1.8cm] (key) at (8.8, 11.0) {\scriptsize Key (病理)};
\node[small_box, fill=drawOrangeFill!50, draw=drawOrangeLine!70,
      minimum width=1.8cm] (val) at (8.8, 10.0) {\scriptsize Value (EHR)};

% Attention heatmap (6x6) inside CMT
\begin{scope}[shift={(11.2, 10.5)}]
  \node[font=\scriptsize\bfseries, text=drawRedLine!70!black] at (1.0, 1.8) {注意力权重};
  \foreach \row in {0,...,5} {
    \foreach \col in {0,...,5} {
      \pgfmathsetmacro{\dist}{abs(\row-\col)}
      \pgfmathtruncatemacro{\noise}{int(mod(\row*3+\col*5+\row*\col,5))}
      \pgfmathtruncatemacro{\score}{min(4, int(\dist + \noise/3))}
      \ifnum\score=0 \def\cellcol{heatDeep}\fi
      \ifnum\score=1 \def\cellcol{heatDark}\fi
      \ifnum\score=2 \def\cellcol{heatMid}\fi
      \ifnum\score=3 \def\cellcol{heatLight}\fi
      \ifnum\score=4 \def\cellcol{heatBg}\fi
      \fill[\cellcol] ({\col*0.32},{-\row*0.28})
        rectangle ({\col*0.32+0.30},{-\row*0.28+0.26});
      \draw[white, very thin] ({\col*0.32},{-\row*0.28})
        rectangle ({\col*0.32+0.30},{-\row*0.28+0.26});
    }
  }
\end{scope}

% MCG module inside CMT
\node[purple_node, minimum width=2.2cm, minimum height=2.2cm,
      inner sep=6pt] (mcg) at (10.5, 7.0)
    {\textbf{MCG}\\[1pt]{\scriptsize 模态门控}\\[-1pt]{\scriptsize CT\,0.45}\\[-1pt]{\scriptsize 病理\,0.35}\\[-1pt]{\scriptsize EHR\,0.20}};

% ================================================================
%  STAGE 3: HIERARCHICAL DIAGNOSIS (x=15~18)
% ================================================================
% Coarse classification
\node[green_node, minimum width=3.2cm, minimum height=1.8cm] (coarse) at (16.5, 13.0)
    {\textbf{粗分类头}\\[1pt]{\footnotesize 5类疾病大类}\\[-1pt]{\scriptsize 准确率 96.2\%}};

% Fine classification (larger, with embedded bar chart)
\node[green_node, minimum width=3.6cm, minimum height=4.5cm] (fine) at (16.5, 8.5) {};
\node[font=\small\bfseries, text=drawGreenLine!80!black] at (16.5, 10.3)
    {细分类头};
\node[font=\footnotesize, text=drawGreenLine!60!black] at (16.5, 9.7)
    {23个具体病种};

% Embedded accuracy comparison bar chart in fine classification box
% Center math: bars 粗(0.3-0.9)+细(1.2-1.8) group center=1.05; zone center
% x=16.5; shift was 15.2 (group abs center 16.25, off-left by 0.25cm).
% Adjusted to 15.45 so group center = 16.5 ✓
\begin{scope}[shift={(15.45, 6.8)}]
  \draw[-{Stealth[scale=0.5]}, drawGreyLine] (0.0, 0.0) -- (0.0, 2.2);
  \draw[-{Stealth[scale=0.5]}, drawGreyLine] (0.0, 0.0) -- (2.8, 0.0);
  % Bars: coarse vs fine accuracy
  \fill[barGreen!70] (0.3, 0.0) rectangle (0.9, 1.92); % 96.2%
  \fill[barBlue!70] (1.2, 0.0) rectangle (1.8, 1.79); % 89.7%
  \node[font=\tiny, text=barGreen!80!black] at (0.6, 2.05) {96.2};
  \node[font=\tiny, text=barBlue!80!black] at (1.5, 1.92) {89.7};
  \node[font=\tiny] at (0.6, -0.2) {粗};
  \node[font=\tiny] at (1.5, -0.2) {细};
  \node[font=\tiny, left, text=drawGreyLine] at (-0.05, 1.0) {50};
  \node[font=\tiny, left, text=drawGreyLine] at (-0.05, 2.0) {100};
  \node[font=\tiny, text=drawGreyLine] at (1.4, -0.7) {\scriptsize 准确率(\%)};
\end{scope}

% Hierarchical arrow between coarse and fine
\draw[data_arrow] (coarse.south) -- (fine.north);
\node[tag, text=drawOrangeLine, fill=white, fill opacity=0.9, text opacity=1,
      anchor=west] at (17.0, 11.4) {\scriptsize 层次化Softmax};

% ================================================================
%  STAGE 4: EXPLAINABILITY (x=20~24)
% ================================================================
% Grad-CAM with mini heatmap
\node[blue_node, minimum width=3.2cm, minimum height=2.5cm] (gradcam) at (21.5, 13.5) {};
\node[font=\footnotesize\bfseries, text=drawBlueLine!80!black] at (21.5, 14.3)
    {Grad-CAM};
% Mini heatmap for Grad-CAM (4x4 grid)
% Fix E.5: was shift=(20.5, 12.6) → bottom row at y=11.55, outside
% node bottom y=12.25 (Grad-CAM node center 13.5, height 2.5).
% Moved up to (20.5, 13.6): rows at 13.6/13.25/12.9/12.55 all inside.
% Center math: 4 cells × 0.4 = 1.6 wide; heatmap center = shift_x + 0.8.
% Zone center x=21.5. shift_x was 20.5 → heatmap center 21.3 (off-left
% by 0.2cm). Adjusted to 20.7 → heatmap center 21.5 ✓
\begin{scope}[shift={(20.7, 13.6)}]
  \foreach \r in {0,...,3} {
    \foreach \c in {0,...,3} {
      \pgfmathtruncatemacro{\v}{int(mod(\r*\r+\c*3+2,5))}
      \ifnum\v=0 \def\cc{heatDeep}\fi
      \ifnum\v=1 \def\cc{heatDark}\fi
      \ifnum\v=2 \def\cc{heatMid}\fi
      \ifnum\v=3 \def\cc{heatLight}\fi
      \ifnum\v=4 \def\cc{heatBg}\fi
      \fill[\cc] ({\c*0.4},{-\r*0.35}) rectangle ({\c*0.4+0.38},{-\r*0.35+0.33});
    }
  }
  \node[font=\tiny, text=drawBlueLine] at (0.8, -1.65) {CT注意力区域};
\end{scope}

% SHAP analysis with horizontal bar chart
\node[orange_node, minimum width=3.2cm, minimum height=3.5cm] (shap) at (21.5, 9.5) {};
\node[font=\footnotesize\bfseries, text=drawOrangeLine!80!black] at (21.5, 10.9)
    {SHAP 分析};
% Embedded horizontal bar chart
\begin{scope}[shift={(20.2, 8.0)}]
  % Feature importance bars
  \fill[barOrange!70] (0.8, 1.6) rectangle (2.3, 1.85); % 主诉
  \fill[barOrange!50] (0.8, 1.15) rectangle (1.9, 1.4); % 检查
  \fill[barOrange!40] (0.8, 0.7) rectangle (1.5, 0.95); % 年龄
  \fill[barOrange!30] (0.8, 0.25) rectangle (1.2, 0.5); % 性别
  % Labels
  \node[font=\tiny, anchor=east] at (0.75, 1.725) {主诉};
  \node[font=\tiny, anchor=east] at (0.75, 1.275) {检查};
  \node[font=\tiny, anchor=east] at (0.75, 0.825) {年龄};
  \node[font=\tiny, anchor=east] at (0.75, 0.375) {性别};
  % Axis
  \draw[drawGreyLine, thin] (0.8, 0.15) -- (0.8, 2.0);
  \node[font=\tiny, text=drawGreyLine] at (1.5, 0.0) {\scriptsize 贡献度};
\end{scope}

% GPT-4 diagnosis report
\node[purple_node, minimum width=3.2cm, minimum height=1.8cm] (gpt4) at (21.5, 5.8)
    {\textbf{GPT-4}\\[1pt]{\footnotesize 诊断报告生成}};

% ================================================================
%  STAGE 5: CLINICAL VALIDATION (x=25~27)
% ================================================================
\node[red_node, minimum width=2.8cm, minimum height=2.5cm,
      inner sep=6pt] (mdt) at (25.5, 10)
    {\textbf{临床验证}\\[2pt]{\scriptsize 双盲审}\\[-1pt]{\scriptsize Kappa=0.87}\\[2pt]{\scriptsize MDT 讨论}};

% ================================================================
%  AUC COMPARISON PANEL (bottom)
% ================================================================
\node[rectangle, rounded corners=6pt, draw=drawBlueLine, fill=drawBlueFill!12,
    minimum width=22cm, minimum height=3.5cm, thick, inner sep=0pt] (aucpanel) at (12, 1.5) {};
% Fix E.3: title was at y=3.0 (anchor=north west) — INSIDE panel top y=3.25
% by 0.25cm, looking clipped. Move to anchor=south west at y=3.35 so the
% title sits clearly ABOVE the panel like a section header.
\node[font=\small\bfseries, text=drawBlueLine!80!black, anchor=south west]
    at (1.5, 3.35) {多模态融合性能对比 (AUC)};

% AUC bar chart
\begin{scope}[shift={(5, 0.3)}]
  % Axes
  \draw[-{Stealth[scale=0.6]}, thick, drawGreyLine] (0.3, 0.0) -- (0.3, 2.3);
  \draw[-{Stealth[scale=0.6]}, thick, drawGreyLine] (0.3, 0.0) -- (10, 0.0);
  % Grid
  \foreach \y/\lab in {0.46/0.5, 0.92/0.6, 1.38/0.7, 1.84/0.8, 2.1/0.9} {
    \draw[gray!20, thin] (0.3, \y) -- (9.5, \y);
    \node[font=\scriptsize, left, text=drawGreyLine] at (0.25, \y) {\lab};
  }
  % CT alone bar
  \fill[barBlue!70, draw=barBlue!90!black, rounded corners=1pt]
      (1.0, 0.0) rectangle (1.8, 1.87); % 0.891
  \node[font=\scriptsize, text=barBlue!80!black] at (1.4, 2.0) {0.891};
  \node[font=\scriptsize] at (1.4, -0.25) {CT};
  % Pathology alone bar
  \fill[barGreen!70, draw=barGreen!90!black, rounded corners=1pt]
      (2.5, 0.0) rectangle (3.3, 1.83); % 0.876
  \node[font=\scriptsize, text=barGreen!80!black] at (2.9, 1.96) {0.876};
  \node[font=\scriptsize] at (2.9, -0.25) {病理};
  % EHR alone bar
  \fill[barOrange!70, draw=barOrange!90!black, rounded corners=1pt]
      (4.0, 0.0) rectangle (4.8, 1.67); % 0.812
  \node[font=\scriptsize, text=barOrange!80!black] at (4.4, 1.8) {0.812};
  \node[font=\scriptsize] at (4.4, -0.25) {EHR};
  % MedFusion bar (emphasized, taller)
  \fill[barRed!60, draw=barRed!80!black, rounded corners=1pt, line width=1pt]
      (5.8, 0.0) rectangle (6.8, 2.0); % 0.943
  \node[font=\scriptsize\bfseries, text=barRed!80!black] at (6.3, 2.13) {0.943};
  \node[font=\scriptsize\bfseries] at (6.3, -0.25) {MedFusion};
  % Improvement annotation — placed further right (8.5 instead of 8.2)
  % so "↑5.8% vs CT" doesn't crowd the "0.943" value label above the bar
  \draw[-{Stealth[scale=0.5]}, very thick, drawRedLine]
      (7.3, 1.4) -- (7.3, 1.9);
  \node[font=\scriptsize\bfseries, text=drawRedLine, anchor=west]
      at (7.4, 1.65) {$\uparrow$5.8\% vs CT};
\end{scope}

% Legend for AUC panel
\begin{scope}[shift={(17, 1.6)}]
  \fill[barBlue!70] (0, 0.6) rectangle (0.3, 0.8);
  \node[font=\scriptsize, anchor=west] at (0.35, 0.7) {CT单模态};
  \fill[barGreen!70] (0, 0.2) rectangle (0.3, 0.4);
  \node[font=\scriptsize, anchor=west] at (0.35, 0.3) {病理单模态};
  \fill[barOrange!70] (2.5, 0.6) rectangle (2.8, 0.8);
  \node[font=\scriptsize, anchor=west] at (2.85, 0.7) {EHR单模态};
  \fill[barRed!60] (2.5, 0.2) rectangle (2.8, 0.4);
  \node[font=\scriptsize, anchor=west] at (2.85, 0.3) {MedFusion};
\end{scope}

% ================================================================
%  CONNECTIONS — Core data flow
% ================================================================
% Feature extraction → Projection
\draw[data_arrow] (ct.east) -- (projCT.west);
\draw[arrow] (path.east) -- (projPath.west);
\draw[arrow] (ehr.east) -- (projEHR.west);

% Projection → CMT
\draw[data_arrow] (projCT.east) -- (cmt.west |- projCT);
\draw[arrow] (projPath.east) -- (cmt.west |- projPath);
\draw[arrow] (projEHR.east) -- (cmt.west |- projEHR);

% CMT → Classification
\draw[data_arrow, rounded corners=6pt]
    (cmt.east |- coarse) -- (coarse.west);
\draw[data_arrow, rounded corners=6pt]
    (cmt.east |- fine) -- (fine.west);

% Classification → Explainability
\draw[arrow, rounded corners=6pt]
    (coarse.east) -- ++(0.8, 0) |- (gradcam.west);
\draw[arrow] (fine.east) -- (shap.west);
\draw[arrow, rounded corners=6pt]
    (fine.east) -- ++(0.8, 0) |- (gpt4.west);

% Explainability → Clinical validation
\draw[arrow, rounded corners=6pt]
    (gradcam.east) -- ++(0.5, 0) |- (mdt.west |- 0,11.5);
\draw[arrow] (shap.east) -- (mdt.west);
\draw[arrow, rounded corners=6pt]
    (gpt4.east) -- ++(0.5, 0) |- (mdt.west |- 0,8.5);

% ================================================================
%  MDT FEEDBACK LOOP (red dashed, U-shaped along bottom)
% ================================================================
% Previous path had a DIAGONAL segment from (25.5, 7.75) to (-1.5, 4.0)
% that cut DIRECTLY ACROSS the entire figure (visually crossing zones,
% boxes, and arrows). Replaced with an orthogonal U-path that goes
% DOWN past the AUC panel bottom (y=-0.25), LEFT to the far margin,
% UP to clients level, then RIGHT into Pathology.
\draw[feed_arrow, rounded corners=8pt, line width=1.2pt]
    (mdt.south) -- (25.5, -1.0) -- (-2.0, -1.0) -- (-2.0, 10) -- (path.west);
% Label moved into the empty space along the bottom horizontal segment
% (was at (12, 3.8) sitting on the old diagonal mid-point).
\node[font=\footnotesize\bfseries, text=drawRedLine,
      fill=white, fill opacity=0.9, text opacity=1,
      inner sep=3pt, rounded corners=2pt] at (12, -0.7) {MDT 反馈优化};

% ================================================================
%  STAGE LABELS (top)
% ================================================================
\node[font=\footnotesize\bfseries, text=drawBlueLine!80!black,
      fill=drawBlueFill!60, draw=drawBlueLine!40, rounded corners=3pt,
      inner sep=4pt] at (1.5, 16.0) {阶段一：多模态特征提取};
\node[font=\footnotesize\bfseries, text=drawGreyLine!80!black,
      fill=drawGreyFill!80, draw=drawGreyLine!40, rounded corners=3pt,
      inner sep=4pt] at (5, 16.0) {投影};
\node[font=\footnotesize\bfseries, text=drawRedLine!80!black,
      fill=drawRedFill!60, draw=drawRedLine!40, rounded corners=3pt,
      inner sep=4pt] at (10.5, 16.0) {阶段二：跨模态注意力融合};
\node[font=\footnotesize\bfseries, text=drawGreenLine!80!black,
      fill=drawGreenFill!60, draw=drawGreenLine!40, rounded corners=3pt,
      inner sep=4pt] at (16.5, 16.0) {阶段三：层次化诊断};
\node[font=\footnotesize\bfseries, text=drawOrangeLine!80!black,
      fill=drawOrangeFill!60, draw=drawOrangeLine!40, rounded corners=3pt,
      inner sep=4pt] at (21.5, 16.0) {阶段四：可解释性};
\node[font=\footnotesize\bfseries, text=drawRedLine!80!black,
      fill=drawRedFill!60, draw=drawRedLine!40, rounded corners=3pt,
      inner sep=4pt] at (25.5, 16.0) {阶段五};

% ================================================================
%  ZONE BACKGROUNDS
% ================================================================
\begin{pgfonlayer}{background}
  % Stage 1: Feature extraction
  \fill[drawBlueFill!10, rounded corners=8pt] (-0.5, 4.5) rectangle (3.5, 15.5);
  \draw[drawBlueLine!30, dashed, thick, rounded corners=8pt] (-0.5, 4.5) rectangle (3.5, 15.5);
  % Projection
  \fill[drawGreyFill!30, rounded corners=8pt] (3.5, 4.5) rectangle (6.5, 15.5);
  \draw[drawGreyLine!30, dashed, thick, rounded corners=8pt] (3.5, 4.5) rectangle (6.5, 15.5);
  % Stage 2: CMT fusion
  \fill[drawRedFill!8, rounded corners=8pt] (6.8, 4.5) rectangle (14.2, 15.5);
  \draw[drawRedLine!25, dashed, thick, rounded corners=8pt] (6.8, 4.5) rectangle (14.2, 15.5);
  % Stage 3: Hierarchical diagnosis
  \fill[drawGreenFill!10, rounded corners=8pt] (14.5, 4.5) rectangle (18.5, 15.5);
  \draw[drawGreenLine!30, dashed, thick, rounded corners=8pt] (14.5, 4.5) rectangle (18.5, 15.5);
  % Stage 4: Explainability
  \fill[drawOrangeFill!10, rounded corners=8pt] (19.5, 4.5) rectangle (23.5, 15.5);
  \draw[drawOrangeLine!30, dashed, thick, rounded corners=8pt] (19.5, 4.5) rectangle (23.5, 15.5);
  % Stage 5: Clinical validation
  \fill[drawRedFill!10, rounded corners=8pt] (23.8, 4.5) rectangle (27.2, 15.5);
  \draw[drawRedLine!25, dashed, thick, rounded corners=8pt] (23.8, 4.5) rectangle (27.2, 15.5);
\end{pgfonlayer}

% ================================================================
%  DATA VOLUME ANNOTATION (centered under 阶段五 label)
% ================================================================
% Was anchor=north east at (27.0, 15.3) — pinned to right edge of Stage 5
% zone. Stage 5 zone spans x=23.8..27.2 (center x=25.5). Changed to
% anchor=north at (25.5, 15.3) so the box is centered horizontally
% under the "阶段五" label (also at center x=25.5).
\node[font=\scriptsize, text=drawGreyLine, fill=drawGreyFill!50,
      draw=drawGreyLine!40, rounded corners=3pt, inner sep=4pt,
      anchor=north] at (25.5, 15.3)
    {\shortstack[l]{3家三甲医院\\12,847例\\多中心数据集}};

\end{tikzpicture}
\end{document}
